\documentclass[11pt]{article}
\usepackage[margin=0.9in]{geometry}
\usepackage{graphicx,amsmath,amssymb,booktabs,hyperref,xcolor}
\graphicspath{{/tmp/cheby_h0/figs/compare/}}
\hypersetup{colorlinks=true,linkcolor=blue,urlcolor=blue}
\setlength{\parskip}{6pt}
\title{Head-to-head: Chebyshev strict-$h{=}0$ vs Linear-CDF kernel-band\\
       on the tabulated-$\mu$ K=3 CRRA REE architecture}
\author{Fixed-Point Factory --- two operators, same architecture, same lookup table}
\date{June 2026}

\begin{document}
\maketitle

\section*{The two operators}

Both share the tabulated-$\mu$ architecture and the same per-cube-cell CRRA
clearing. They differ only in how $P(u_1,u_2,u_3)$ is represented and how
$\mu(p, u_k)$ is built:

\begin{center}
\begin{tabular}{lll}
\toprule
& \textbf{Chebyshev (strict $h{=}0$)} & \textbf{Linear-CDF (kernel band)} \\
\midrule
$u$-grid                 & Lobatto $\times$ atanh ($c{=}2$), $|u|\!\leq\!9.9$
                         & Gaussian quantiles, $|u|\!\leq\!2.33$ \\
$P$ representation       & tensor-Chebyshev polynomial
                         & values + bilinear interp \\
slice eval               & exact polynomial value
                         & bilinear interp on raw values \\
contour integral         & strict $h{=}0$ POU + chebroots
                         & Gaussian band $K_h$, $h=0.5$ \\
GL quadrature order      & $N_Q$ auto-selected per $\gamma$
                         & $N_Q=16$ fixed \\
Lookup table $\mu(p,u_k)$ & built at $G_p=121$ p-points
                         & same \\
Per-cell clearing         & CRRA bisection on the lookup
                         & same \\
\bottomrule
\end{tabular}
\end{center}

\noindent
The Chebyshev approach computes the \emph{exact} strict-$h{=}0$ co-area
integral using polynomial root extraction; the linear-CDF approach uses
a \emph{regularised} integral with finite kernel bandwidth.

\section{Wall-time per $\Phi$ and convergence to machine $\varepsilon$}

\begin{center}
\begin{tabular}{lcc}
\toprule
operator & ms / $\Phi$ at $G{=}7$ & cold $\rightarrow$ machine-$\varepsilon$ solve \\
\midrule
Chebyshev strict $h{=}0$, $N_Q{=}16$ & 540 & 110 s (43s Anderson + 67s NK) \\
Linear-CDF kernel band $h{=}0.5$    &   6 &   0.3 s (Anderson alone) \\
\bottomrule
\end{tabular}
\end{center}

\textbf{Linear-CDF kernel band is $\sim 300\times$ faster.}
And reaches machine $\varepsilon$ at every $\gamma$ in the sweep on the first
try (no NQ tuning needed because the operator is analytic in $P$).

\begin{center}
\includegraphics[width=0.9\linewidth]{03_floors_compare.png}
\end{center}

The Chebyshev solver has 3 of 9 $\gamma$ values stalling at $\sim 10^{-3}$
even with NQ auto-search (sweet-spot dependence on $\gamma$). The Linear-CDF
kernel band has no such issue --- all $\gamma$ nail to machine $\varepsilon$.

\section{Equilibrium statistics vs $\gamma$}

\begin{center}
\includegraphics[width=0.9\linewidth]{01_slope_compare.png}
\end{center}

\begin{center}
\begin{tabular}{ccccccc}
\toprule
& \multicolumn{2}{c}{slope $\alpha^*$} & \multicolumn{2}{c}{$1-R^2_{\rm lin}$}
& \multicolumn{2}{c}{$1-R^2_{\rm nonparam}$} \\
$\gamma$ & Cheb & Lin-CDF & Cheb & Lin-CDF & Cheb & Lin-CDF \\
\midrule
0.10 & 0.3182 & 0.2054 & 0.198 & 0.253 & 0.0361 & 0.0331 \\
0.30 & 0.3039 & 0.2282 & 0.167 & 0.129 & 0.0339 & 0.0104 \\
0.50 & 0.3179 & 0.2558 & 0.175 & 0.063 & 0.0309 & 0.0024 \\
0.80 & 0.3292 & 0.2837 & 0.158 & 0.025 & 0.0306 & 0.0000 \\
1.00 & 0.3447 & 0.2957 & 0.145 & 0.016 & 0.0285 & 0.0001 \\
1.50 & 0.3675 & 0.3141 & 0.111 & 0.008 & 0.0175 & 0.0004 \\
2.00 & 0.3826 & 0.3245 & 0.101 & 0.005 & 0.0129 & 0.0006 \\
3.00 & 0.4150 & 0.3359 & 0.090 & 0.002 & 0.0076 & 0.0005 \\
5.00 & 0.4493 & 0.3459 & 0.082 & 0.001 & 0.0033 & 0.0003 \\
\bottomrule
\end{tabular}
\end{center}

\paragraph{Two distinct FPs.} The strict-$h{=}0$ Chebyshev and the
kernel-band linear-CDF give \emph{different} fixed points (different
slope $\alpha^*$ at every $\gamma$). The kernel-band is an
$h{>}0$-regularised version of the strict integral and converges to a
different operator's FP.

\begin{center}
\includegraphics[width=0.95\linewidth]{02_deficit_compare.png}
\end{center}

\textbf{Both} show the same qualitative trend: slope rises with $\gamma$,
deficit falls. But the kernel-band's nonparametric deficit is essentially
zero from $\gamma\!\geq\!0.5$ --- \emph{the kernel-band FP is almost
exactly a function of $T$ alone}, with no one-to-one breakdown.

The strict-$h{=}0$ Chebyshev FP keeps a nonzero one-to-one breakdown
(0.03 at $\gamma{=}1$, falling to 0.003 at $\gamma{=}5$).

\paragraph{Why the difference matters.}
\begin{itemize}
  \item The kernel-band integral $\int K_h(P-p)\,f f\, du$ averages out
        fine-scale variation of the co-area integrand. As $h \to 0$, it
        converges to the strict integral; but at $h=0.5$ there is
        substantial smoothing.
  \item This smoothing has the effect of making $P$ closer to a function
        of $T$ alone. The kernel-band operator is essentially solving a
        \emph{less informationally rich} REE: it's the limit where the
        market price summarises only the aggregate signal.
  \item The strict-$h{=}0$ Chebyshev FP captures genuine non-$T$
        information that the kernel-band loses.
\end{itemize}

\section{Price contours $\{P(u_1, u_2, u_3) = p\}$}

\subsection{$\gamma=1$ comparison side-by-side}

\begin{center}
\includegraphics[width=\linewidth]{04_contours_compare.png}
\end{center}

The Chebyshev contours (left, on the wider Lobatto-atanh grid) bend more
than the Linear-CDF contours (right, on the narrower Gaussian-quantile grid).
Both show the expected diagonal structure (the contour $P=p$ runs roughly along
$u_1+u_2=\,$const), but the curvatures differ.

\subsection{Chebyshev contours across $\gamma$}

\begin{center}
\includegraphics[width=\linewidth]{06_cheb_contours_all.png}
\end{center}

\subsection{Linear-CDF contours across $\gamma$}

\begin{center}
\includegraphics[width=\linewidth]{05_lin_cdf_contours_all.png}
\end{center}

\section{logit-$P$ vs $T$ scatter at $\gamma=1$}

\begin{center}
\includegraphics[width=0.95\linewidth]{07_scatter_compare_g1.png}
\end{center}

The Chebyshev scatter has visible spread around any vertical line
($T = $ const): genuine one-to-one breakdown. The Linear-CDF scatter
looks like a smooth single curve --- essentially a function of $T$.

\section{Summary}

\begin{itemize}
  \item \textbf{Chebyshev strict $h{=}0$ + POU}: exact contour integral,
        the actual K=3 strict-$h{=}0$ REE. Slow ($\sim 540$ ms/$\Phi$,
        $\sim 110$ s/solve), $N_Q$ needs auto-tuning per $\gamma$, 6 of 9
        $\gamma$ reach machine $\varepsilon$. Captures nonzero one-to-one
        breakdown (0.003--0.036).
  \item \textbf{Linear-CDF kernel band, $h=0.5$}: regularised integral,
        smooth operator. Fast ($\sim 6$ ms/$\Phi$, $\sim 0.3$ s/solve), all
        9 $\gamma$ reach machine $\varepsilon$ on first try. But the FP
        is a smoothed version --- nonparametric deficit collapses to $\sim 0$
        for $\gamma\!\geq\!0.5$.
  \item \textbf{Implication for the paper}: the one-to-one $P \leftrightarrow T$
        breakdown is a feature of the \emph{strict} $h{=}0$ operator,
        present in the Chebyshev path but smoothed away by kernel
        regularisation. To report the genuine breakdown, use Chebyshev
        strict-$h{=}0$ + POU; to get fast smooth FPs for parametric
        sweeps (with the caveat of smoothing), use Linear-CDF kernel.
\end{itemize}

\end{document}
